\chapter{Overall Conclusion and Future Scope}

\section{Conclusion}
This project designed, implemented and rigorously evaluated machine-learning pipelines across all five units of the
course. Every result was measured on held-out data.
\begin{itemize}
\item \textbf{Unit 1:} cleaning and exploratory analysis turned two messy market datasets into quantified insights.
Free apps dominate installs (92.2\% of apps; median 100,000 vs 1,000 installs for paid apps). Video-game sales are
concentrated in hits (the top 1\% of titles earn 21.1\% of sales) and peaked in 2008.
\item \textbf{Unit 2:} tuned XGBoost predicted power-plant output with RMSE 3.17 MW ($R^2$ 0.966), 31\% lower error than
linear regression. An imbalance-aware, leakage-free XGBoost model identified bank-deposit subscribers with F1 0.531 and
a 3.3$\times$ lift over random calling.
\item \textbf{Unit 3:} K-Means and Gaussian Mixture clustering recovered the static/dynamic structure of smartphone
activity data (accuracy up to 0.612 without labels), and showed that geometric clusters differ from semantic activities.
\item \textbf{Unit 4:} the forward–backward and Viterbi algorithms were implemented from first principles and verified
against \texttt{hmmlearn} (difference $<10^{-9}$). A supervised hybrid HMM decoded the activity sequences of unseen
people with 98.7\% accuracy; its Markov chain removed 63\% of the per-window errors.
\item \textbf{Unit 5:} ten ensemble classifiers were implemented and compared with Logistic Regression and Decision Tree
baselines for diabetes prediction. The best model, \dbBestName{} with an out-of-fold-tuned threshold, reached
\textbf{F1 \dbBestFone{} and ROC-AUC \dbBestAuc{}} (accuracy \dbBestAccPct) on \dbTest{} unseen patients, against F1
\dbLrFone{} for Logistic Regression and \dbDtFone{} for the Decision Tree. The improvement is statistically significant,
and SHAP explanations agree with clinical guidelines.
\end{itemize}
The central lesson is that good machine learning depends as much on \emph{methodology} as on algorithms. The methodology
used throughout was: careful cleaning, leakage control, metrics matched to the class balance, cross-validated model
selection, thresholds chosen on training data, and interpretation. Within that methodology, combining learners, whether
over models (ensembles) or over time (HMMs), consistently gave the most accurate and reliable predictions.

\section{Future Scope}
\begin{enumerate}
\item \textbf{External and prospective validation:} test the diabetes model on an independent hospital cohort with raw,
continuous laboratory values, and the HAR model on free-living data with different phone positions.
\item \textbf{Cost-sensitive decisions:} replace F1-optimal thresholds with thresholds derived from real costs (e.g.\ the
cost of a missed diabetic vs an unnecessary test, or of a call vs a deposit).
\item \textbf{Broader tuning:} Bayesian optimisation (e.g.\ Optuna) for all leading models, and weighted or
learned-weight voting ensembles.
\item \textbf{Richer sequential models:} semi-supervised HMMs that need only a few labels, hidden semi-Markov models with
explicit activity durations, and neural sequence models for comparison.
\item \textbf{Deployment and interpretability:} package the diabetes model as a web-based risk calculator that shows
per-patient SHAP explanations to clinicians, with monitoring for data drift.
\item \textbf{Fairness:} audit model performance across gender and age groups before any clinical use.
\end{enumerate}
